\documentclass{cumcmthesis}
% 电子版提交时通常可使用：\documentclass[withoutpreface]{cumcmthesis}

\usepackage{booktabs}
\usepackage{amsmath,amssymb,bm}
\usepackage{subcaption}
\usepackage{graphicx}
\usepackage{float}
\usepackage{algorithm}
\usepackage{algorithmic}
\usepackage{cleveref}
% cumcmthesis v2.6 自定义了旧版 \supercite；在加载 biblatex 前释放该命令，
% 由 biblatex 提供统一的上标引用实现。
\let\supercite\relax
\usepackage[backend=biber,style=numeric,sorting=none]{biblatex}
\addbibresource{references.bib}

% 按当届题面和队伍信息修改，不在模板中固定年份、题号和小问数。
\newcommand{\HSKProblemCode}{A}
\newcommand{\HSKTeamNumber}{0000}
\newcommand{\HSKContestYear}{YYYY}
\newcommand{\HSKContestMonth}{9}
\newcommand{\HSKContestDay}{1}

\title{基于XXX模型的XXX问题研究}
\tihao{\HSKProblemCode}
\baominghao{\HSKTeamNumber}
\schoolname{学校名称}
\membera{成员A}
\memberb{成员B}
\memberc{成员C}
\supervisor{}
\yearinput{\HSKContestYear}
\monthinput{\HSKContestMonth}
\dayinput{\HSKContestDay}

\begin{document}
\maketitle

\begin{abstract}
% 按实际小问数量写作。每问采用“模型—算法—关键数值—结论”结构，删除未使用的占位语句。
本文围绕XXX问题，针对题目给出的XXX数据、XXX机制与XXX约束，建立XXX模型并完成求解与检验。

针对问题一，……。

% 继续按实际小问添加摘要段，不固定为三问。

本文通过参数扰动、数据噪声、替代模型或多算法验证中的适用方法检验结论稳定性，并给出模型适用边界。

\keywords{关键词1\quad 关键词2\quad 关键词3\quad 关键词4}
\end{abstract}

\section{问题重述}
\subsection{问题背景}
% 用自己的语言概括现实背景、研究对象和业务机制，不照抄题干。

\subsection{问题要求}
% 概括各小问；内部题目覆盖表不得进入终稿正文或附录。

\section{问题分析}
\subsection{问题一分析}
% 说明直接目标、隐含目标、数据依赖、题型、核心难点和方法选择理由。

% 按实际小问继续添加分析小节。

\section{模型假设与符号说明}
\subsection{模型假设}
\begin{enumerate}
    \item 假设一：……。说明设立原因、影响、失效偏差和检验方式。
    \item 假设二：……。
    \item 假设三：……。
\end{enumerate}

\subsection{符号说明}
\begin{table}[H]
    \centering
    \caption{主要符号说明}
    \label{tab:symbols}
    \begin{tabular}{cccc}
        \toprule
        符号 & 含义 & 单位 & 类型 \\
        \midrule
        $x_i$ & 决策变量示例 & -- & 决策变量 \\
        $s_t$ & 状态变量示例 & -- & 状态变量 \\
        $\theta$ & 参数示例 & -- & 参数 \\
        \bottomrule
    \end{tabular}
\end{table}

\section{数据预处理}
\subsection{数据口径与字段说明}
% 说明数据表、字段、单位、时间/空间粒度、样本范围和关联键。

\subsection{数据质量与变换}
% 说明缺失、异常、重复、单位换算、时间对齐及特征构造依据。

\section{模型建立与求解}
\subsection{问题一模型建立与求解}
\subsubsection{模型建立}
设……，目标函数为
\begin{equation}
    \min Z = C_1 + C_2 + C_3,
    \label{eq:q1-objective}
\end{equation}
其中，$C_1$、$C_2$ 和 $C_3$ 分别表示……。

\subsubsection{求解方法}
% 写求解器、参数、终止条件、程序位置和可行性检查。

\subsubsection{结果与分析}
% 图表必须就近解释，并绑定标准结果工作簿和 MATLAB 绘图脚本。

% 按实际小问继续添加“模型建立—求解方法—结果与分析”。

\section{模型检验}
% 按题型选择有效性检验、残差诊断、边界验证、多算法对比或误差分析。

\section{敏感性与鲁棒性分析}
% 参数扰动、数据噪声、场景压力测试或替代模型对比。

\section{模型评价与改进}
\subsection{模型优势}
\subsection{模型局限}
\subsection{改进方向}

\section{结论}
% 逐问直接回答题目，给出关键数值、策略和边界，不写空泛总结。

\printbibliography[title={参考文献}]

\begin{appendices}
\section{核心代码与复现说明}
% 正文只保留伪代码；完整代码可放附录或独立附件，并说明输入、功能和输出。

\section{补充图表}
% 仅放正文不宜展开但仍有证据价值的图表。
\end{appendices}

\end{document}
